\chapter{Retorno nonádico, memoria y holonomía de Cartan}
\label{chap:puente-retorno-holonomia-cartan}
\index{retorno!27--54--108}\index{holonomía!TPK--Cartan}
\index{memoria!distinción respecto de torsión}

La realización gravitatoria no comienza identificando una variable discreta
con una magnitud geométrica. Comienza por distinguir tres niveles: el retorno
del transporte observable, la memoria que diferencia historias con igual
proyección y la holonomía de una conexión. El puente que se formula aquí
recibe los retornos y cociclos ya demostrados en HMT; no los reconstruye
desde la acción de Holst, desde una masa ni desde una constante metrológica.

\section{Jerarquía exacta de los retornos}
\label{sec:jerarquia-retornos-27-54-108}

Sea \(F_{\rm obs}\) la cronología observable definida en el núcleo HMT y sea
\[
 \mathcal K_{27}:=F_{\rm obs}^{27}.
\]
El símbolo \(\mathcal K_{27}\) designa exclusivamente la composición de
veintisiete actualizaciones del lazo posicional y orientado; no designa una
ventana genérica, una traza de \(108\) pasos ni el estado TPK enriquecido
completo. Denotemos por \(p\) la proyección a las posiciones de los dos
cursores y por \(\mathsf J_{\rm or}\) la inversión simultánea de sus
orientaciones. La demostración primaria se encuentra en
\cref{prop:retorno-fobs-identidad}; sus consecuencias tipadas son
\[
 p(\mathcal K_{27}x)=p(x),\qquad
 o(\mathcal K_{27}x)=\mathsf J_{\rm or}\,o(x),
 \qquad \mathsf J_{\rm or}^{\,2}=I,
\]
\[
 F_{\rm obs}^{54}:
 \text{restaura posiciones y orientaciones},
 \qquad
 F_{\rm obs}^{108}=I:
 \text{restaura además la fase observable}.
\]
Por tanto, retorno de posición, retorno de orientación y retorno de fase son
afirmaciones diferentes. Ninguna de ellas implica por sí sola el retorno de
todos los campos de una extensión enriquecida.

El cociente de modos y los cociclos de orientación se establecen en
\cref{sec:cociclos-modo-orientacion,thm:descenso-necesario-fav}. En el orden
areal--volumétrico,
\[
 N_{27}=9F_{\rm av},
 \qquad
 F_{\rm av}=
 \begin{pmatrix}0&1\\1&1\end{pmatrix}.
\]
Si
\[
 \Omega_{\rm sw}(\gamma)=\sum_{a\subset\gamma}\omega_{\rm sw}(a),
 \qquad
 S(\gamma)=\sum_{a\subset\gamma}|\omega_{\rm sw}(a)|,
 \qquad
 G_{\rm or}(\gamma)=\sum_{a\subset\gamma}g(a),
\]
entonces un bloque de veintisiete pasos contiene tres vueltas nonádicas y una
inversión de orientación:
\[
 \Omega_{\rm sw}(\mathcal K_{27})=0,\qquad
 S(\mathcal K_{27})=18,\qquad
 G_{\rm or}(\mathcal K_{27})=1.
\]
En el retorno observable de \(108\) pasos,
\[
 \Omega_{\rm sw}(C_{108})=0,\qquad
 S(C_{108})=72,\qquad
 G_{\rm or}(C_{108})=4.
\]
Así, para el funcional aditivo declarado
\[
 \mathcal A_\ast(\gamma)
 :=G_{\rm or}(\gamma)+\lambda_\ast S(\gamma),
\]
se obtienen, sin introducir una escala dimensional,
\[
 \mathcal A_\ast(\mathcal K_{27})=1+18\lambda_\ast,
 \qquad
 \mathcal A_\ast(C_{108})=4+72\lambda_\ast.
\]
Estos valores son conteos del registro. Su posterior uso térmico o
gravitatorio no altera su procedencia.

\subsection{Sistema dirigido de coberturas, refinamiento y límite de la conexión de Cartan}
\label{sec:cobertura-acumulativa-dirigida}

Sea \(\mathcal O_{108}\) la órbita observable y, para cada arista cronológica
\(e\), sea
\[
 c(e)=\bigl(g(e),s(e)\bigr)\in\mathbb N^2
\]
el par formado por los indicadores de inversión de orientación y de cambio
efectivo de hoja. Para una historia dirigida
\(\gamma=e_1\cdots e_n\), se define
\[
 c(\gamma)=\sum_{j=1}^{n}c(e_j).
\]
La ley de concatenación toma la forma
\[
 c_{m+n}(x)
 =c_m(x)+c_n\!\left(F_{\rm obs}^{m}x\right).
\]
Los conteos anteriores expresan, de manera uniforme sobre la órbita,
\[
 c_{27}(x)=(1,18),
 \qquad
 c_{108}(x)=(4,72).
\]

\begin{theorem}[Cobertura acumulativa de la cronología observable]
\label{thm:cobertura-acumulativa-dirigida}
Sobre
\[
 \widetilde{\mathcal O}_{108}^{+}
 :=\mathcal O_{108}\times\mathbb N^2
\]
defínase
\[
 \widetilde F^{+}(x,m)
 :=\bigl(F_{\rm obs}x,m+c(x\to F_{\rm obs}x)\bigr).
\]
Entonces el transporte de veintisiete pasos satisface
\[
 \widetilde{\mathcal K}_{27}^{+}(x,m)
 =\bigl(\mathcal K_{27}x,m+(1,18)\bigr),
\]
y
\[
 \left(\widetilde{\mathcal K}_{27}^{+}\right)^4(x,m)
 =\bigl(x,m+(4,72)\bigr)\neq(x,m).
\]
Por tanto, su proyección observable tiene orden cuatro, mientras que la
acción acumulativa del monoide cronológico posee orden semigrupal infinito.
\end{theorem}

\begin{proof}
La ley cocíclica y la constancia de \(c_{27}\) dan
\[
 c_{108}(x)
 =\sum_{j=0}^{3}c_{27}(\mathcal K_{27}^{j}x)
 =4(1,18)=(4,72).
\]
La primera coordenada retorna por
\(\mathcal K_{27}^{4}=I\); la segunda no retorna. Tras \(q>0\) iteraciones,
el incremento es \(q(4,72)\), distinto de cero.
\end{proof}

\begin{remark}[Tipo de la extensión acumulativa]
La construcción anterior amplía historias dirigidas mediante contadores
positivos. No es un levantamiento automorfo del grupoide
\(\mathcal G_{\rm mem}\), ni demuestra que la coordenada ilimitada pertenezca
ya al estado \(\mathcal X_{\mathrm{TPK}}^{\mathrm{enr}}\). Su completación de Grothendieck
\(\mathbb Z^2\) puede utilizarse para análisis armónico, pero no transforma
el coste positivo
\(\mathcal A_\ast=G_{\rm or}+\lambda_\ast S\) en un carácter del grupo de
isotropía: al invertir formalmente una ruta, sus indicadores no adquieren
signo opuesto.
\end{remark}

\paragraph{Dominio de validez de la extensión acumulativa.}
La jerarquía \(27\to54\to108\), las incidencias de \(F_{\rm av}\) y los
cociclos anteriores son identidades exactas en la cronología y los
lectores declarados. La afirmación se restringe al estado observable que
esos lectores conservan; extenderla al estado enriquecido completo requeriría
probar el retorno de cada campo olvidado.

\section{Grupoide de memoria}
\label{sec:grupoide-memoria-tpk}

Sea \(\mathsf P_{\rm TPK}^{\rm enr}\) el grupoide de caminos enriquecidos:
sus objetos son estados TPK supervivientes y sus flechas son rutas admisibles
con origen y destino declarados, composición por concatenación e inversión
por reversión de ruta. Sea \(\sim_{\rm mem}\) una equivalencia que sólo
identifica dos flechas cuando conserva los campos de memoria requeridos por
el lector considerado y que es congruente con identidades, inversión y
concatenación. Definimos
\[
 \mathcal G_{\rm mem}
 :=\mathsf P_{\rm TPK}^{\rm enr}/\!\sim_{\rm mem}.
\]
Esta definición no transforma la memoria en torsión. La memoria pertenece al
estado discreto y distingue historias; la torsión pertenece a una conexión
geométrica y mide \(D_\omega e\). Sólo un mapa entre ambos dominios puede
relacionarlas.

Una realización geométrica exige, como datos adicionales, una longitud
\(\ell_0>0\), una representación del grupoide
\[
 \rho_{\rm C}:
 \mathcal G_{\rm mem}
 \longrightarrow
 \operatorname{Spin}(1,3)\ltimes\RR^{1,3},
\]
y un mapa de campos
\[
 \mathfrak R_{\rm grav}:
 \mathsf P_{\rm TPK}^{\rm enr}
 \dashrightarrow
 \mathsf{Cartan}_{1,3}.
\]
Para definir el segundo mapa deben especificarse su acción sobre objetos y
rutas, su compatibilidad con concatenación e inversión y la regularidad de la
conexión resultante. Ninguna de estas propiedades se deduce del mero conteo
de retornos.

\section{Conexión conjunta y condición de entrelazamiento}
\label{sec:conexion-cartan-conjunta}

Bajo las hipótesis geométricas del \cref{chap:gravedad} ---variedad
orientada de dimensión cuatro, co-tétrada \(e^I\), conexión de Lorentz
\(\omega^{IJ}\) y condiciones de borde allí declaradas---, la conexión de
Cartan conjunta se escribe
\[
 \mathcal A_{\ell_0}
 =
 \begin{pmatrix}
  \omega&\ell_0^{-1}e\\
  0&0
 \end{pmatrix}.
\]
Su curvatura es la identidad algebraica
\[
 \mathcal F_{\ell_0}
 =\dd\mathcal A_{\ell_0}
  +\mathcal A_{\ell_0}\wedge\mathcal A_{\ell_0}
 =
 \begin{pmatrix}
  F_\omega&\ell_0^{-1}T_{\rm tor}\\
  0&0
 \end{pmatrix},
\]
donde
\[
 F_\omega^{IJ}
 =\dd\omega^{IJ}+\omega^I{}_K\wedge\omega^{KJ},
 \qquad
 T_{\rm tor}^{I}=D_\omega e^I.
\]
Curvatura y torsión aparecen como componentes distintas de un mismo defecto
de transporte. La igualdad matricial se sigue algebraicamente de los datos
explícitos \(e,\omega,\ell_0\), pero no construye estas entradas desde TPK.

El contenido fuerte del puente es la conmutatividad de holonomías:
\[
 \boxed{
 \Hol_{\mathcal A_{\ell_0}}
 \bigl(\mathfrak R_{\rm grav}(\gamma)\bigr)
 =
 \rho_{\rm C}\!\left(
 \Hol_{\rm TPK}(\gamma)
 \right)}
 \qquad
 (\gamma\in\mathsf P_{\rm TPK}^{\rm enr}).
\]
La igualdad debe ser compatible con identidades, inversión, concatenación,
subdivisión, cambio de punto base y conjugación. Sólo entonces la memoria
discreta queda representada en una conexión; no queda rebautizada como
torsión ni como curvatura.

\begin{criterion}[Condiciones de existencia del puente TPK--Cartan]
\label{crit:puente-tpk-cartan}
No existe una realización con la conmutatividad y la compatibilidad declaradas
si ocurre cualquiera de los siguientes hechos:
\begin{enumerate}
 \item dos representantes equivalentes de una misma ruta producen
 holonomías geométricas no conjugadas;
 \item la imagen no preserva identidades, inversos o concatenación;
 \item una subdivisión admisible altera la holonomía;
 \item la igualdad de holonomías falla para \(\mathcal K_{27}\) o para sus
 concatenaciones de profundidades \(54\) y \(108\);
 \item el mapa identifica memoria y torsión sin exhibir un morfismo que
 preserve sus tipos;
 \item alguna elección descendente ---la acción de Holst, \(G\), una masa o
 una carta SI--- se utiliza para redefinir \(A\), \(C^\ast\) o el retorno
 TPK de origen.
\end{enumerate}
\end{criterion}

\paragraph{Antecedencia del funcional \(\Gamma_{6\to8}\) respecto de la realización Palatini--Cartan--Holst.}
El \cref{mdg:chap:barbero} fija, una vez construidos \(A,C^\ast,q_\pm\), el
funcional interno \(\Gamma_{6\to8}\). El \cref{chap:gravedad} estudia después
la acción Palatini--Cartan--Holst y sus ecuaciones variacionales. El puente
presente no introduce una ruta inversa entre ambas etapas.
